% SPDX-FileCopyrightText: 2026 Will Cook
% SPDX-License-Identifier: CC-BY-4.0
\documentclass[11pt]{article}
\input{problem-note-preamble}
\renewcommand{\commit}{99f4bf47422abbd8757cbb22b50ba079d764d3a7}
\renewcommand{\commitshort}{99f4bf47422a}
\hypersetup{
  pdftitle={A basis for the 2-kernel of Euler's totient},
  pdfsubject={Erdős Problem Notes: Problem 249},
  pdfkeywords={irrationality, Euler totient, k-regular sequences, 2-kernel, Lambert series, Möbius inversion, linear independence, Lean 4},
}

\newcommand{\Stot}{S}
\newcommand{\EB}{E}
\newcommand{\Aw}{A}
\newcommand{\Ker}{\mathcal K}
\newtheorem{equivform}[theorem]{Equivalent formulation}

% Status tag from the claim taxonomy, printed on every numbered result.
\newcommand{\stag}[1]{{\normalfont\footnotesize\textsc{[#1]}}}

\runninghead{Erd\H{o}s \#249: a basis for the dyadic totient sections}
\title{A Basis for the 2-Kernel of Euler's Totient}
\subtitle{Exact finite-level ranks and consequences for Erd\H{o}s
Problem~\#249}
\author{Will Cook}
\date{July 2026}

\begin{document}
\maketitle
\noteprovenance

\begin{problem}[first-harmonic anti-concentration on supplier fibres]
\label{prob:firstharmonic}
For each $h\ge1$, do there exist $s\ge1$ and $0<\eta<1$ such that, for every
$X_0$, there are $X,L$ with
\[
 \max(X_0,1)\le X,\qquad h\le L-s,\qquad
 16(2X+h+L+2)\le2^L,
\]
for which the four bounds below hold?
\end{problem}

\begin{abstract}
For every $e\ge1$, the dyadic sections of Euler's totient through level $e$
have rational rank $2^e+1$. The two zero-residue sections and the odd-residue
sections form a basis, and two families of scalar reductions generate every relation.
Divisibility makes an evaluation matrix diagonal modulo a prime, with nonzero
diagonal entries. The construction gives rank $k^e+1$ in every base $k\ge2$
and an integral normal form.
Separately, $\sum_{n\ge1}(\ph(n)\bmod m)2^{-n}$ is irrational for every
$m\ge3$, with a complete rationality classification at dyadic moduli.
For $S=\sum_{n\ge1}\ph(n)2^{-n}$, one nonintegral tail shift forces
full-depth certificates in every sufficiently late pair of multipliers.
Irrationality is equivalent to an explicit central-residue gap for every
prospective denominator; its arithmetic supply remains open.
\end{abstract}



\section{A basis and all its relations}\label{sec:results}
Set $\ph(0)=0$. For $k\ge2$, the $k$-kernel consists of the sequences
$n\mapsto\ph(k^jn+r)$ with $j\ge0$ and $0\le r<k^j$.
The following theorem gives a unique coordinate system at every finite level.
\begin{theorem}[exact rank of every totient $k$-kernel truncation]
\label{thm:kkernelrank}
Let $k\ge2$ and $e\ge1$ be integers, write $F^{(k)}_{j,r}(n)=\ph(k^jn+r)$, and put
\[
 V_{k,e}=\operatorname{span}_{\Q}
 \{\,F^{(k)}_{j,r}:0\le j\le e,\ 0\le r<k^j\,\}.
\]
Then $\dim_{\Q}V_{k,e}=k^e+1$, and
\[
 \mathcal B_{k,e}
 =\{F^{(k)}_{0,0},F^{(k)}_{1,0}\}
 \cup\{\,F^{(k)}_{j,r}:1\le j\le e,\ 1\le r<k^j,\ k\nmid r\,\}
\]
is a basis.  Every omitted section reduces to a basis element by an explicit
scalar: $F^{(k)}_{j,0}=k^{j-1}F^{(k)}_{1,0}$ for $j\ge1$, and if $r=k^tu$ with
$t=\max\{s:k^s\mid r\}\ge1$ and $k\nmid u$, then
\[
 F^{(k)}_{j,r}=C_k(t,u)\,F^{(k)}_{j-t,u},
 \qquad
 C_k(t,u)=k^t\prod_{\substack{p\mid k\\ p\nmid u}}\Bigl(1-\tfrac1p\Bigr).
\]
\end{theorem}
\begin{enumerate}
\item\label{res:rank} \stag{unconditional progress}
\textbf{Exact finite-level rank.}  Theorem~\ref{thm:kkernelrank} is the
coordinate statement: every totient $k$-kernel truncation has rank $k^e+1$.
\end{enumerate}
The retained condition is $k\nmid r$. At $k=6$, for example, the residue
$r=2$ is retained. At base $2$ and level $3$, the fifteen sections have nine
coordinates; the omitted positive-residue sections satisfy
\[
 \ph(8n+2)=\ph(4n+1),\qquad
 \ph(8n+4)=2\ph(2n+1),\qquad
 \ph(8n+6)=\ph(4n+3).
\]

\subsection{The evaluation matrix}\label{sec:rank}
The two ingredients are scalar reduction and arithmetic independence.
The local formula for $\ph$ gives the displayed reductions. For a prime
$p\mid k$ that does not divide $u$, the argument $k^{j-t}n+u$ is a $p$-adic
unit, so multiplication by $k^t$ contributes its power of $p$ and one factor
$1-1/p$. When $p\mid u$, the argument already contains $p$, and only its
power contributes. This gives $C_k(t,u)$.

For independence, consider finitely many forms $L_i(n)=a_in+b_i$ with
positive slopes and $a_ib_j\ne a_jb_i$ for $i\ne j$. A finite shift makes
the intercepts positive. Write $L_i=g_iA_i$, where $A_i$ is primitive.
Choose an odd prime $\ell$ avoiding the slopes, all $g_i\ph(g_i)$, and the
nonzero cross determinants. In row $i$, choose distinct primes
$q_{ij}\equiv1\pmod\ell$ for $j\ne i$, outside the same finite exceptional
set, and impose
\[
 L_j(n)\equiv0\pmod{q_{ij}},\qquad A_i(n)\equiv2\pmod\ell.
\]
The Chinese remainder theorem combines these congruences. The resulting
progression for $A_i(n)$ is reduced: a prime $q_{ij}$ dividing $A_i(n)$
would also divide the nonzero cross determinant; $\ell$ is excluded by the
chosen residue; and primitivity excludes the primes of the slope.
Dirichlet's theorem supplies arbitrarily large primes $p=A_i(n)$ with
$p\nmid g_i$. Only the selected affine form is made prime in this row.
The diagonal value is
$\ph(g_i)(p-1)\not\equiv0\pmod\ell$. Each off-diagonal value is divisible
by $\ell$, since $q_{ij}-1$ divides the corresponding totient.
One evaluation per row gives
\[
 \bigl(\ph(L_j(n_i))\bigr)_{i,j}
 \equiv\operatorname{diag}(u_1,\ldots,u_s)\pmod\ell,
 \qquad u_i\ne0.
\]
The determinant is nonzero over $\Q$.

Apply this construction after restricting a relation among the retained
sections to $n=km+1$. The positive-residue sections give pairwise
nonproportional forms
\[
 k^{j+1}m+(k^j+r),\qquad k\nmid r,
\]
and the remaining form is $km+1$. Equality between two affine fractions
would make the residue at the larger level divisible by $k$. The matrix
therefore eliminates every positive-residue coefficient. If $A,B$ are the
two zero-residue coefficients, the restriction gives
$A+B\ph(k)=0$. Evaluation at $n=k$ gives $A+Bk=0$, since
$\ph(k^2)=k\ph(k)$. As $\ph(k)<k$, both coefficients vanish.
The reductions give spanning and the count is
\[
 2+\sum_{j=1}^e(k^j-k^{j-1})=k^e+1.
\]
This proves Theorem~\ref{thm:kkernelrank}. The scalar reduction at level
$j-t$, with the Euler product written out, is
\href{https://github.com/wcook04/plectis-erdos/blob/f36a98bf3d3e6f65f1074e3b800e3293d5b8a51a/ErdosProblems/Erdos249/PaperCompleteR7/KernelIntegral.lean#L226}{Lean source}.

\begin{corollary}[Dyadic basis]\label{res:basis}
The family
\[
 \{\ph_{0,0},\ph_{1,0}\}
 \cup\{\ph_{j,r}:j\ge1,\ 0<r<2^j,\ r\text{ odd}\},
 \qquad \ph_{j,r}(n)=\ph(2^jn+r),
\]
is a basis for the rational span of the full dyadic kernel. Every rational
relation is generated by the scalar reductions. The complete level-zero
truncation has dimension one.
\end{corollary}
\begin{proof}
Every finite subfamily lies in a sufficiently high truncation. The theorem
gives its independence; the reductions give equality of the full spans.
The level-zero clause is
\href{https://github.com/wcook04/plectis-erdos/blob/f36a98bf3d3e6f65f1074e3b800e3293d5b8a51a/ErdosProblems/Erdos249/PaperCompleteR7/ArithmeticAssemblies.lean#L112}{Lean source}.
\end{proof}

\begin{corollary}[Integral normal form]\label{cor:integral-normal-form}
Let $k\ge2$ and $e\ge1$. The canonical family is a $\Z$-basis of the
module generated by the sections through level $e$. For each omitted
channel $F_i$, write its canonical reduction as $F_i=a_iF_{j(i)}$, where
$j(i)$ is retained and $a_i$ is a nonnegative integer. In the free abelian
group on all the channels, the vectors
\[
 R_i=E_i-a_iE_{j(i)}\qquad(i\text{ omitted})
\]
form a $\Z$-basis of the kernel of evaluation $E_i\mapsto F_i$.
In particular, every integral relation has a unique integral expression
in these elementary relations, and their rank is
$\sum_{j=1}^{e-1}k^j$.
\end{corollary}
\begin{proof}
The denominator $\prod_{p\mid k,\,p\nmid u}p$ in $C_k(t,u)$ divides $k$,
so the successive channel reductions have integer scalars. They give
integral coordinates in the retained family; its rational independence
makes these coordinates unique over $\Z$ as well.

Now let $x=\sum_i x_iE_i$ evaluate to zero. For each omitted index $i$,
subtract $x_iR_i$. The result is supported on retained indices and still
evaluates to zero, so independence forces it to vanish. Thus the $R_i$
generate the integral relation module. Conversely, the coefficient at
an omitted index $i$ in $\sum_h c_hR_h$ is exactly $c_i$: each relation
has coefficient one at its own omitted coordinate and zero at all the
others. This proves independence and uniqueness without dividing by
any integer. There are
\[
 (1+k+\cdots+k^e)-(k^e+1)=\sum_{j=1}^{e-1}k^j
\]
omitted channels. The formal statement combines the integral coordinate
basis, these two-term relation vectors, and their exact rank.
This complete integral normal form is \href{https://github.com/wcook04/plectis-erdos/blob/25ef6245d15a47548c6926369ae8f1a0f0a14a80/ErdosProblems/Erdos249/PaperCompleteR8/KernelRelationBasis.lean#L398}{kernel-checked in Lean}.
\end{proof}

For example, at base $2$ and depth $2$, retain
$F_{0,0},F_{1,0},F_{1,1},F_{2,1},F_{2,3}$. The two omitted channels satisfy
$F_{2,0}=2F_{1,0}$ and $F_{2,2}=F_{1,1}$, so
\[
 E_{2,0}-2E_{1,0},\qquad E_{2,2}-E_{1,1}
\]
are an integral basis of all relations among the seven channels. The
rank calculation alone would give two; the basis also identifies every
relation and its unique coefficients.

The module here is the one generated by the channels. It is not saturated
in the group of all integer-valued sequences: at base $2$ and $e\ge2$,
$\ph(4n+3)/2$ is integer-valued but has canonical coefficient $1/2$.

\subsection{A reusable consequence}
\begin{corollary}[Freezing periodic coefficients]\label{cor:periodic-freezing}
If $L_i$ are pairwise nonproportional positive-slope affine forms and $w_i$
are rational-valued periodic sequences, then
\[
 \sum_i w_i(n)\ph(L_i(n))=0\quad\text{for all sufficiently large }n
 \quad\Longrightarrow\quad w_i(n)=0\quad\text{for every }i,n.
\]
\end{corollary}
\begin{proof}
Choose a common period $Q$. On each progression $n=a+Qt$, the coefficients
are constant and the cross determinants become $Q(a_ib_j-a_jb_i)\ne0$.
Shift beyond the threshold and apply the affine independence just proved.
The statement is \href{https://github.com/wcook04/plectis-erdos/blob/f36a98bf3d3e6f65f1074e3b800e3293d5b8a51a/ErdosProblems/Erdos249/PaperCompleteR7/PeriodicAndPulse.lean#L64}{Lean source},
with a separate period for each coefficient.
\end{proof}
The same argument preserves the basis after deletion of a finite prefix.
For the two zero channels, large $n\equiv1\pmod k$ and large powers of $k$
replace the two evaluations used above.
\section{Bounded residues and rationality}\label{sec:family}
The residue theorem uses the same separation principle: force neighbouring
coefficients to vanish while retaining one central value.

\begin{theorem}\label{res:residueseries}
For every $m\ge3$,
\[
 A_m=\sum_{n\ge1}\frac{\ph(n)\bmod m}{2^n}\notin\Q.
\]
For $k\ge1$ and $f:\Z/2^k\Z\to\Q$, the series
$\sum_{n\ge1}f(\ph(n)\bmod2^k)2^{-n}$ is rational exactly when $f$ is
constant on the even residue classes. If that constant is $c$, its value is
$3f(1)/4+c/4$.
\end{theorem}

\begin{lemma}[A bounded isolated pulse]\label{lem:bounded-pulse}
Let $a_n\in\Z$ satisfy $|a_n|\le C$. Suppose that for arbitrarily large $L$
there is $N>L$ such that $a_N\ne0$ and
$a_{N+t}=0$ for $0<|t|\le L$. Then $\sum_{n\ge1}a_n2^{-n}$ is irrational.
\end{lemma}
\begin{proof}
If the sum has denominator $q$, every scaled tail
$T_j=\sum_{t\ge1}a_{j+t}2^{-t}$ belongs to $q^{-1}\Z$.
Choose $L$ so large that $C2^{-L}<1/q$.
The right zero block gives $|T_N|\le C2^{-L}<1/q$, hence $T_N=0$.
The left zero block now gives
$T_{N-L-1}=a_N2^{-L-1}$, a nonzero element of $q^{-1}\Z$ with absolute
value less than $1/q$, a contradiction. The statement is
\href{https://github.com/wcook04/plectis-erdos/blob/f36a98bf3d3e6f65f1074e3b800e3293d5b8a51a/ErdosProblems/Erdos249/PaperCompleteR7/PeriodicAndPulse.lean#L106}{Lean source}.
\end{proof}

\begin{proof}[Proof of Theorem~\ref{res:residueseries}]
Fix a unit $s\pmod m$ and a block length $L$. For each nonzero
$t\in[-L,L]$, choose a distinct prime $q_t\equiv1\pmod m$ larger than $L$
and outside the prime divisors of $m$. Prescribe
\[
 p\equiv s\pmod m,\qquad p\equiv-t\pmod{q_t}.
\]
This is a reduced progression, because each $q_t>|t|$. Dirichlet supplies
arbitrarily large primes $p$ in it. Then $m\mid\ph(p+t)$ for every
$0<|t|\le L$, whereas $\ph(p)\equiv s-1\pmod m$.

Subtract $f(0)$ from the observable and clear its denominators. Whenever
$f(s-1)\ne f(0)$, the preceding construction gives arbitrarily long bounded
isolated pulses, so Lemma~\ref{lem:bounded-pulse} applies. For the
least-residue map choose $s=-1$: its central residue is $m-2\ne0$ when
$m\ge3$. For $m=2^k$, every even residue $r$ has $r+1$ a unit, so a
nonconstant restriction to the even residues supplies such a pulse.
Conversely, $\ph(n)$ is even for every $n\ge3$. Constancy on the even
classes therefore makes all terms after $n=2$ constant, giving the stated
rational value. This converse requires no recurrent-support theorem.
The three clauses together are
\href{https://github.com/wcook04/plectis-erdos/blob/f36a98bf3d3e6f65f1074e3b800e3293d5b8a51a/ErdosProblems/Erdos249/PaperCompleteR7/RationalObservableClassification.lean#L277}{Lean source}.
\end{proof}
The bounded-pulse proof fixes its coefficient bound before choosing the block.
For $S$, the termwise condition $2^L\mid\ph(N+L)$, with $N+L>1$, forces
$2^L\le\ph(N+L)\le N+L-1$. Hence the tail envelope
$(N+L+2)2^{-L}>1$ cannot exclude any positive denominator.
The following criterion instead measures cancellation in a weighted prefix.

\section{The actual series and the remaining residue gap}\label{sec:carry-rank}
Write
\[
 S=\sum_{n\ge1}\ph(n)2^{-n},\qquad
 R_N=\sum_{j\ge1}\ph(N+j)2^{-j},\qquad
 \Delta_h(N)=R_{N+h}-R_N.
\]
The prefix identity $2^NS=\Phi_N+R_N$, with
$\Phi_N=\sum_{n=1}^N\ph(n)2^{N-n}\in\Z$, gives
\begin{equation}\label{eq:tail-phase}
 \Delta_h(N)\equiv2^N(2^h-1)S\pmod\Z.
\end{equation}
In particular a nonintegral tail difference is a precise arithmetic
statement about $S$.

\subsection{One seed removes the full-depth restriction}
For $h,N,L\in\N$, let
\[
 D_{h,N,L}=\sum_{j=0}^{L-1}
   \bigl(\ph(N+h+1+j)-\ph(N+1+j)\bigr)2^{L-1-j},
 \qquad B=N+h+L+2.
\]
Write $K(h,N,L)$ for
\[
 B<D_{h,N,L}\bmod2^L<2^L-B.
\]
The exact truncation and $\ph(n)\le n$ give
\begin{equation}\label{eq:tail-error}
 |2^L\Delta_h(N)-D_{h,N,L}|\le B.
\end{equation}
Thus some $L$ satisfies $K(h,N,L)$ exactly when $\Delta_h(N)\notin\Z$:
necessity follows from the strict inequalities; sufficiency follows by
choosing $L$ with $2^L\|\Delta_h(N)\|_{\R/\Z}>2B$.

\begin{theorem}[Full-depth amplification]\label{res:fulldepth}
Fix $d\ge1$ and $N\ge0$. If $\Delta_d(N)\notin\Z$, then every sufficiently
late pair $\{t,t+1\}$ contains an $m$ such that $K(md,N,md)$ holds.
Consequently
\[
 \exists t\ge1:\ K(td,N,td)
 \quad\Longleftrightarrow\quad\Delta_d(N)\notin\Z.
\]
Requiring this for every $d\ge1,N\ge0$ is equivalent to $S\notin\Q$.
\end{theorem}
\begin{proof}
Set $F_t=\Delta_{td}(N)$. Equation~\eqref{eq:tail-phase} yields
$F_{t+1}\equiv2^dF_t+F_1\pmod\Z$. If $K(td,N,td)$ fails,
\eqref{eq:tail-error} gives
\[
 \|F_t\|_{\R/\Z}\le\epsilon_t,
 \qquad \epsilon_t=2(N+2td+2)2^{-td}.
\]
Two adjacent failures would imply
$\|F_1\|\le2^d\epsilon_t+\epsilon_{t+1}$, impossible for all sufficiently
large $t$ because the left side is positive and the right side tends to zero.
Conversely, if $F_1$ is integral, then every $F_t$ is integral by the same
recurrence, so no certificate exists. Finally, \eqref{eq:tail-phase} is
nonintegral for every $d,N$ when $S$ is irrational. If
$S=a/(2^cv)$ with $v$ odd, take $N=c$ and $d=\ph(v)$ to make it integral.
\end{proof}
For instance $D_{1,12,16}=-143140$, with residue $53468$ modulo $2^{16}$
and $B=31$, supplies a seed. The amplifier produces infinitely many
full-depth certificates from this finite fact; it does not provide seeds
for all $d,N$.

\subsection{A single canonical endpoint}\label{sec:frontier}
For $c\ge0$, odd $v\ge1$ and a positive multiple $H$ of $\ph(v)$, put
\[
 M=\frac{2^H-1}{v},\qquad
 B_{H,c}=\sum_{j=0}^{H-1}\ph(c+1+j)2^{H-1-j},\qquad K=c+H+1.
\]
Euler's theorem makes $M$ integral.

\begin{equivform}[Canonical Mersenne gap]\label{res:canonicalmersenne}
The series $S$ is irrational if and only if, for every $c\ge0$ and odd
$v\ge1$, some positive multiple $H$ of $\ph(v)$ satisfies
\begin{equation}\label{eq:canonical-gap}
 K<(-B_{H,c})\bmod M<M-K.
\end{equation}
\end{equivform}
\begin{proof}
The identity
\[
 B_{H,c}=2^HR_c-R_{c+H}=M(vR_c)-\Delta_H(c)
\]
is exact. For $H\ge1$, $|\Delta_H(c)|<c+H+1=K$. Indeed $R_N\le N+1$
for $N\ge1$ and $0<R_0\le3/2$; the strict inequalities follow by subtracting
a positive tail. If $S=a/(2^cv)$, then $vR_c\in\Z$ and
$\Delta_H(c)\in\Z$. The residue of $-B_{H,c}$ is therefore within distance
less than $K$ of $0\pmod M$, so \eqref{eq:canonical-gap} fails.

If $S$ is irrational, fix $c,v$. Then $vR_c$ is nonintegral. As $H$ tends
to infinity through positive multiples of $\ph(v)$,
$\Delta_H(c)/M\to0$ and $K/M\to0$. Hence the fractional part of
$-B_{H,c}/M=-vR_c+\Delta_H(c)/M$ stays away from both endpoints of $[0,1]$,
which gives \eqref{eq:canonical-gap} for all sufficiently large such $H$.
\end{proof}
This proof identifies the equivalent missing quantifier. It does not supply
$H$ arithmetically without assuming the desired irrationality. The
equivalence, with the modulus written as the natural quotient $(2^H-1)/v$
and $v$ odd, is
\href{https://github.com/wcook04/plectis-erdos/blob/f36a98bf3d3e6f65f1074e3b800e3293d5b8a51a/ErdosProblems/Erdos249/PaperCompleteR7/ArithmeticAssemblies.lean#L156}{Lean source}.
The finite residues also satisfy
\[
 \rho_{H,N+1,M}=
 \bigl(2\rho_{H,N,M}-\ph(N+H+1)+\ph(N+1)\bigr)\bmod M,
 \quad \rho_{H,N,M}=(-B_{H,N})\bmod M,
\]
which is an exact recurrence for searches directed at \eqref{eq:canonical-gap}.

\section{Information that does not force the gap}\label{sec:nogo}
\begin{enumerate}
\item\label{res:actualorbit} \stag{exact reformulation}
\textbf{Actual LCM orbit frontier.}  Let
$H_a=\lcm(1,\ldots,2^a)$ and
$R_a=\operatorname{totientTail}(2H_a)-\operatorname{totientTail}(H_a)$.
Then $\Stot$ is irrational if and only if, beyond every threshold $a_0$,
some $a\ge a_0$ has $R_a\notin\Z$.  This is an exact cofinal
nonintegrality criterion for the actual LCM diagonal, not a quantitative
separation estimate and not a proof that such indices occur.
{\footnotesize\emph{Checked:}
\mword{Erdos249257/TotientActualLcmOrbitNonintegrality.lean}{37}{irrational_totientSeries_iff_actualLcmOrbitNonintegralitySupply}{exact frontier},
\mword{Erdos249257/TotientActualLcmOrbitNonintegrality.lean}{53}{irrational_totientSeries_of_actualLcmOrbitNonintegralitySupply}{forward implication}.}

\item\label{res:actualsign} \stag{unconditional progress}
\textbf{Positive LCM corridor and top-edge boundary.}  For $a\ge8$ and
$J+(a+6)<2\cdot2^a$, the true translated orbit difference
$\operatorname{totientTail}(2H_a+J)-\operatorname{totientTail}(H_a+J)$ is
strictly positive.  If it is represented by an integer and the dyadic
modulus has room for the directed endpoint strip, the carry survivor is
strictly negative and the discrepancy residue is exactly the top-edge
representative $2^K-\mathrm{carry}_K$, lying strictly between the strip's lower
edge and $2^K$.  Thus the positive corridor eliminates the nonnegative true
survivor but does not itself produce a central-band contradiction; an
independent arithmetic exclusion of the top boundary is still required.
{\footnotesize\emph{Checked:}
\mword{Erdos249257/TotientActualLcmOrbitSign.lean}{39}{actualLcmTailDiff_shift_pos}{positive corridor},
\mword{Erdos249257/TotientActualLcmOrbitSign.lean}{172}{actualLcm_trueEndpointSurvivor_neg}{negative survivor},
\mword{Erdos249257/TotientActualLcmOrbitSign.lean}{211}{actualLcm_integral_forces_topEdgeResidue}{top-edge residue}.}

\item\label{res:factorideal} \stag{unconditional progress}
\textbf{LCM factor-ideal shift-algebra no-go.}  For every $t\ge3$, with
$H=\lcm(1,\ldots,t)$, there is a nonzero synthetic dyadic coboundary whose
forcing letters reproduce the actual totient differences at $t-2$ prescribed
indices and lie in the ideal generated by $\varphi(H)$.  Every finite integer
shift polynomial preserves the exact coboundary cancellation, the $H$-factor
ideal and all lower factor ideals, together with uniform state and letter
bounds.  A sparse-anchor variant also preserves every prescribed whole-ray
anchor.  Hence factor ideals, finite shift-algebra closure and these bounds
cannot by themselves force a contradiction.  The witness is synthetic: its
letters are not claimed to be actual totient differences, nonlinear
combinations are outside the theorem, and no cofinal family of certificates
follows.
{\footnotesize\emph{Checked:}
\mword{Erdos249257/LcmFactorIdealPulseObstruction.lean}{798}{lcm_factorIdeal_finiteRank_shiftAlgebra_not_sufficient}{finite shift-algebra no-go}.}
\end{enumerate}

The long record supplies a rational control with
$c(n)=\ph(n)$ for odd $n$, $|c(n)-\ph(n)|\le2$ for even $n$, $0\le c(n)\le n$,
and
\[
 \sum_{n\ge1}c(n)2^{-n}=\frac54.
\]
It has an integral tempered carry with dyadic rank at least $2^e-1$ at every
level. Thus the same lower bound for a hypothetical totient carry is
compatible with rationality. For one fixed prime $p$, the exact laws
$g(pn)=p g(n)$ when $p\mid n$ and $g(pn)=(p-1)g(n)$ otherwise, together
with $g-\ph=o(n)$, force $g=\ph$: a nonzero defect keeps a nonzero ratio
to the argument along $p^j n$. The control fails these laws; its full
coefficient rank is asserted only at measured finite depths.

A different obstruction concerns rational approximation. Put
\[
 \Theta_2=\sum_{d\ge1}\frac{\mu(d)}{(2^d-1)^2}=S-\frac12,
 \qquad
 Q(e,Y)=\frac{\bigl(\sum_{d=1}^{Y}\mu(d)/(2^d-1)^{e+2}\bigr)^2}
 {\sum_{d=1}^{Y}\mu(d)/(2^d-1)^{2e+2}}.
\]
The divisor-convolution identity gives the displayed value of $\Theta_2$,
as recorded by Fan~\cite{fan2026totient}. The following exact extremum is
stated in the long record as Theorem~\texttt{thm:rankonefloor}; its proof is
\lword{Erdos249/RankOneSharpFloor.lean}{534}{rankOneSubrankQuotient_eq_one_five_iff}{the sharp-floor argument}
in the linked source.

\begin{theorem}[Positive rank-one floor]\label{res:rankonefloor}
For $e\ge1$ and $Y\ge4$, the denominator of $Q(e,Y)$ is positive, and the
unique minimiser is $(e,Y)=(1,5)$. Every admissible quotient and every
nonempty finite positive weighted average of such quotients satisfies
\[
 Q-\Theta_2>\frac{21}{320}.
\]
The uniform bound $Q-\Theta_2>1/15$ is false already at $(1,5)$.
\end{theorem}
This theorem treats averages after the quotients are formed. Signed
coefficients, coupling before the quotient, and higher-rank kernels require
separate arguments. For any proposed rational approximants $p_j/q_j$ in
lowest terms, a sufficient arithmetic target is
\[
 0<q_j\left|\Theta_2-p_j/q_j\right|\longrightarrow0.
\]
Ordinary convergence and Hankel nonvanishing do not contain this denominator
estimate. The canonical gap \eqref{eq:canonical-gap} remains the equivalent
endpoint; the approximation estimate is one stronger sufficient route.

\section{Antecedents and proof record}\label{sec:open}
Allouche and Shallit introduced $k$-regular sequences~\cite{allouche-shallit}.
Coons proved that $\ph$ is not $k$-regular~\cite{coons}; Martin's stronger
affine ordering theorem implies the independence in
Section~\ref{sec:rank}~\cite{martin-phi-inequalities}. The basis theorem
gives the exact finite-level normal form. Wong~\cite{wong2015} proved
least-residue irrationality when the denominator base equals the modulus;
Section~\ref{sec:family} fixes base two and classifies dyadic observables.
The original question appears in Erd\H{o}s and Graham~\cite[p.~61]{erdosgraham1980}.

The all-base arithmetic reduction and spanning, and exact rank conditional on
the explicit affine-section linear-independence hypothesis, have proofs in the
supplied Lean source snapshot.  The residue-series irrationality, full-depth
amplification, canonical equivalence and positive rank-one floor are also
Lean-checked:
\lword{Erdos249/FullDepthRayAmplifier.lean}{364}{eventually_twoSyndetic_fullDepthKillMultipliers_of_seed}{amplification},
\lword{Erdos249/CyclotomicAnchoredKill.lean}{3165}{fullMersenneCanonicalBasepointResidueGapSupply_iff_irrational}{canonical gap}, and
\lword{Erdos249/RankOneSharpFloor.lean}{624}{rankOneSubrankQuotient_sub_theta_two_gt_twentyOne_div_threeTwenty}{rank-one floor}.
The historical source links retain \texttt{\commitshort}; the all-base
module requires a matching public source supplement.  The unconditional
all-base basis and rank theorem is an ordinary conclusion above: it discharges
the linear-independence hypothesis using Martin's external theorem.  The
integral normal form also has its ordinary proof above. The periodic transport
lemma is checked in \texttt{PeriodicTotientIndependence.lean}; the eventual
version follows by the displayed threshold shift. The rational control and
the termwise-window and prime-dilation boundaries are proved in
\texttt{ParityPerturbedRationalControl.lean} and
\texttt{PrefixValuationAndControlRigidity.lean}, respectively.
Finite evaluation ranks and the all-level carry lower bound remain distinct.
The long record retains the separate continued-fraction bounds
$q\ge2^{39989}$ and $q>10^{12038}$ under rationality, alongside the
Lean-checked Farey exclusion. Neither implies the universal gap.
\appendix
\section{Guide to the formal sources}\label{app:sources}

The public snapshot used by this note is recorded by the following source links.

\mword{Erdos249257/TotientMahlerDefect.lean}{160}{totientKernel_zero_residue}{checked},
\mword{Erdos249257/TotientMahlerDefect.lean}{169}{totientKernel_even_residue_reduce}{checked},
\mword{Erdos249257/TotientTailPeriodKiller.lean}{150}{two_pow_mul_totient_series_eq}{checked},
\href{https://github.com/wcook04/plectis-erdos/blob/99f4bf47422abbd8757cbb22b50ba079d764d3a7/Erdos249257/TotientTailPeriodKiller.lean#L72}{definition},
\mword{Erdos249257/LcmConeFlatness.lean}{316}{exists_certifiedKill_iff_tail_diff_notMem_int}{checked},
\mword{Erdos249257/TotientTailPeriodKiller.lean}{404}{certifiedKill_all_small}{checked},
\mword{Erdos249257/TotientTailPeriodKiller.lean}{407}{tail_diff_not_int_small}{checked},
\mword{Erdos249257/GenericTailOrbitRigidity.lean}{67}{IsTemperedBinaryOrbit}{definition},
\lword{Erdos249/CyclotomicAnchoredKill.lean}{2322}{fullMersenneBlockResidue_succ}{residue recurrence},
\lword{Erdos249/CyclotomicAnchoredKill.lean}{3132}{fullMersenneCenteredResidueGapSupply_of_canonicalBasepoint}{dominance over the remote condition},
\lword{Erdos249/CyclotomicAnchoredKill.lean}{3165}{fullMersenneCanonicalBasepointResidueGapSupply_iff_irrational}{exact equivalence},
\lword{Erdos249/FullDepthRayAmplifier.lean}{364}{eventually_twoSyndetic_fullDepthKillMultipliers_of_seed}{eventual two-syndeticity},
\lword{Erdos249/FullDepthRayAmplifier.lean}{388}{exists_fullDepthKill_on_ray_iff_shift_notMem_int}{pointwise equivalence},
\lword{Erdos249/FullDepthRayAmplifier.lean}{406}{apFullDepthEscape_iff_irrational}{global equivalence},
\lword{Erdos249/FullDepthRayAmplifier.lean}{421}{cofinalFullDepthKillSupply_iff_periodMultipleKillSupply}{cofinal depth deletion},
\lword{Erdos249/FullDepthRayAmplifier.lean}{438}{cofinalFullDepthKillSupply_iff_irrational}{cofinal endpoint},
\mword{Erdos249257/TotientMahlerDefect.lean}{1265}{linearIndependent_oddCoreTotientKernelFamily}{independence},
\mword{Erdos249257/TotientMahlerDefect.lean}{1380}{span_range_fullTotientKernel_eq_span_range_oddCore}{span equality},
\mword{Erdos249257/TotientMahlerDefect.lean}{1392}{totientDyadicSectionBasis}{basis},
\mword{Erdos249257/TotientMahlerDefect.lean}{935}{linearIndependent_canonicalTotientKernelFamily}{independence},
\mword{Erdos249257/TotientMahlerDefect.lean}{989}{finrank_canonicalTotientKernel_eq}{dimension},
\mword{Erdos249257/TotientMahlerDefect.lean}{1084}{finrank_totientKernelThroughLevelFamily_eq}{the exact finite-truncation rank},
\mword{Erdos249257/TotientMahlerDefect.lean}{1145}{not_finiteDimensional_span_fullTotientKernel}{infinite-dimensionality directly},
\lword{Erdos249/RankOneSharpFloor.lean}{534}{rankOneSubrankQuotient_eq_one_five_iff}{unique minimizer},
\lword{Erdos249/RankOneSharpFloor.lean}{624}{rankOneSubrankQuotient_sub_theta_two_gt_twentyOne_div_threeTwenty}{uniform $21/320$ floor},
\lword{Erdos249/RankOneSharpFloor.lean}{634}{rankOneSubrankQuotient_sub_theta_two_gt_one_div_sixteen}{$1/16$ floor},
\lword{Erdos249/RankOneSharpFloor.lean}{645}{not_forall_rankOneSubrankQuotient_sub_theta_two_gt_one_div_fifteen}{$1/15$ failure},
\lword{Erdos249/RankOneSharpFloor.lean}{655}{positive_direct_sum_sub_theta_two_gt_twentyOne_div_threeTwenty}{positive direct sums},
\lword{Erdos249/RankOneSharpFloor.lean}{693}{primitive_form_abs_gt_twentyOne_div_threeTwenty}{rational linear-form obstruction},
\mword{Erdos249257/TotientKernelIndex.lean}{124}{totientKernelSectionIndex_spec}{coordinate specification},
\mword{Erdos249257/TotientKernelIndex.lean}{213}{existsUnique_totientKernelResidueAtLevel}{unique representation},
\mword{Erdos249257/TotientKernelIndex.lean}{241}{card_totientKernelSectionLevel}{level count},
\mword{Erdos249257/TotientKernelIndex.lean}{277}{card_totientKernelIndex}{full index count},
\mword{Erdos249257/CertificateKernel.lean}{18371}{tsum_totient_div_pow_two_ne_int_div_of_den_le_79639646646701375323355774875831053}{exclusion},
\mword{Erdos249257/CertificateKernel.lean}{18384}{tsum_totient_div_pow_two_ne_ratCast_of_den_le_79639646646701375323355774875831053}{reduced-denominator form},
\mword{Erdos249257/GapFareyBound.lean}{225}{gap_check_window_1_240_first_failure}{first failure},
\mword{Erdos249257/TotientCarryKernelRigidity.lean}{300}{not_irrational_totientSeries_implies_unbounded_carryRank_unconditional}{the barrier},
\mword{Erdos249257/TotientCarryKernelRigidity.lean}{211}{finrank_canonicalCarryKernel_ge_of_linearIndependent}{rank transport from \ref{res:rank}},
\mword{Erdos249257/TotientTailCarryPeriod.lean}{224}{not_irrational_totientSeries_implies_mod_period_and_unbounded_rank}{the modular-period/rank frontier},
\mword{Erdos249257/TotientTailCarryPeriod.lean}{126}{totient_forcing_vanishes_mod_of_dvd_multiplier}{forcing vanishing},
\mword{Erdos249257/TotientTailCarryPeriod.lean}{140}{totient_carry_modEq_geometric_of_dvd_multiplier}{geometric reduction},
\mword{Erdos249257/CertificateKernel.lean}{18445}{tsum_primWeight_div_two_pow_sub_one_eq_totient_series}{identity},
\mword{Erdos249257/MersenneLambertLadder.lean}{370}{primWeight_nonneg}{sign},
\mword{Erdos249257/MersenneLambertLadder.lean}{297}{primWeight_apply_prime}{primes},
\mword{Erdos249257/MersenneLambertLadder.lean}{272}{primWeight_apply_prime_pow}{prime powers},
\mword{Erdos249257/MersenneLambertLadder.lean}{321}{primWeight_not_bounded}{unboundedness},
\mword{Erdos249257/GcdMomentCalculus.lean}{216}{tsum_one_div_mersenne_sq_eq_sigma_sub_tau_series}{divisor-count layer},
\mword{Erdos249257/GcdMomentCalculus.lean}{235}{tsum_totient_div_mersenne_sq_eq_gcd_moment_series}{first gcd moment},
\mword{Erdos249257/SquaredMersenneDiagonalEnclosure.lean}{138}{scaleDiagonalTailDifference_sub_lambertProjectedDiagonal}{exact centre decomposition},
\mword{Erdos249257/SquaredMersenneDiagonalEnclosure.lean}{408}{abs_scaleDiagonalTailDifference_sub_lambertProjectedDiagonal_le}{sharp enclosure},
\mword{Erdos249257/SquaredMersenneDiagonalEnclosure.lean}{426}{scaleFullTarget_miss_of_lambert_projected_separation}{separation criterion},
\mword{Erdos249257/FullTargetPrimeAdjunctionNoGo.lean}{139}{scaleFullTargetHit_iff_integral}{target equivalence},
\mword{Erdos249257/FreshPrimeDeficitDecomposition.lean}{91}{endpointFreshDeficit_nonneg}{endpoint nonnegativity},
\mword{Erdos249257/FreshPrimeDeficitDecomposition.lean}{184}{finiteForeignChannelIncrement_eq_endpointFreshDeficit_sub}{foreign-channel identity},
\mword{Erdos249257/FreshPrimeDeficitDecomposition.lean}{269}{oldFivePoint_sub_secondBranchAdverseDeficit_le_actual}{five-point loss bound},
\mword{Erdos249257/SquareCRTCube.lean}{297}{exists_squareCRT_clean_totient_family}{simultaneous finite family},
\mword{Erdos249257/SquareCRTCube.lean}{327}{exists_squareCRT_clean_horizon}{finite horizon},
\mword{Erdos249257/SquareCRTCube.lean}{459}{squareCRT_clean_block_can_vanish}{vanishing block},
\mword{Erdos249257/SquareCRTCube.lean}{477}{squareCRT_clean_block_can_be_nonzero}{nonzero block},
\mword{Erdos249257/DiagonalFreshLossBridge.lean}{2667}{canonicalAdjacentSuffixCentral_iff_slack_nonneg}{centrality equivalence},
\mword{Erdos249257/DiagonalFreshLossBridge.lean}{2757}{canonicalAdjacentSuffixJumpCentralSupply_iff_slackSupply}{equivalence of the two conditions},
\mword{Erdos249257/DiagonalFreshLossBridge.lean}{2932}{irrational_totientSeries_of_canonicalAdjacentSuffixJumpSlackSupply}{conditional irrationality theorem},
\mword{Erdos249257/ActualForeignResidueProjection.lean}{276}{abs_foreignTailWindow_le_foreignComplementBound}{finite complement bound},
\mword{Erdos249257/ActualForeignResidueProjection.lean}{308}{finiteResidueDiagonal_eq_projectedForeign_add_divisor}{exact channel split},
\mword{Erdos249257/ActualForeignResidueProjection.lean}{414}{scaleFullTarget_miss_of_projected_separation}{conditional target-miss theorem},
\mword{Erdos249257/TotientActualLcmOrbitNonintegrality.lean}{37}{irrational_totientSeries_iff_actualLcmOrbitNonintegralitySupply}{exact frontier},
\mword{Erdos249257/TotientActualLcmOrbitNonintegrality.lean}{53}{irrational_totientSeries_of_actualLcmOrbitNonintegralitySupply}{forward implication},
\mword{Erdos249257/TotientActualLcmOrbitSign.lean}{39}{actualLcmTailDiff_shift_pos}{positive corridor},
\mword{Erdos249257/TotientActualLcmOrbitSign.lean}{172}{actualLcm_trueEndpointSurvivor_neg}{negative survivor},
\mword{Erdos249257/TotientActualLcmOrbitSign.lean}{211}{actualLcm_integral_forces_topEdgeResidue}{top-edge residue},
\mword{Erdos249257/LcmFactorIdealPulseObstruction.lean}{798}{lcm_factorIdeal_finiteRank_shiftAlgebra_not_sufficient}{finite shift-algebra no-go},
\mword{Erdos249257/LcmFactorIdealPulseObstruction.lean}{866}{lcm_factorIdeal_sparseAnchor_not_sufficient}{sparse-anchor no-go},
\mword{Erdos249257/CertificateKernel.lean}{12811}{irrational_ratWeightSeries_eventuallyPeriodic}{the periodic theorem},
\mword{Erdos249257/PrimitiveDeterminantLift.lean}{169}{no_fixed_integral_primitive_euler_index}{no fixed index},
\mword{Erdos249257/PrimitiveDeterminantLift.lean}{148}{twoTierPrimorial_dvd_of_integral_jet}{the primorial divisibility},
\mword{Erdos249257/LcmConeFlatness.lean}{399}{irrational_totient_series_iff_pointwise_certificates}{pointwise},
\mword{Erdos249257/LcmConeFlatness.lean}{412}{irrational_totient_series_iff_certificate_supply}{cofinal at each shift},
\mword{Erdos249257/LcmConeFlatness.lean}{426}{irrational_totient_series_iff_lcm_diagonal_certificate_supply}{$\lcm$ diagonal},
\mword{Erdos249257/PivotAntiReconstruction.lean}{1765}{dtwWindowSeparatedPairs_iff_irrational_totient_series}{window-separated pairs},
\mword{Erdos249257/CertificateKernel.lean}{16025}{irrational_tsum_totient_div_pow_two_of_gap_certificate_supply}{the conditional implication},
\mword{Erdos249257/CertificateKernel.lean}{8335}{irrational_erdosBorwein_series}{checked},
\mword{Erdos249257/CertificateKernel.lean}{18434}{tsum_moebius_div_two_pow_sub_one_eq_half}{checked},
\mword{Erdos249257/CertificateKernel.lean}{18429}{tsum_totient_div_two_pow_sub_one_eq_two}{checked},
\mword{Erdos249257/CertificateKernel.lean}{18454}{totient_series_eq_half_add_moebius_mersenne_square}{checked},
\mword{Erdos249257/CertificateKernel.lean}{18557}{totient_series_eq_half_add_visible_coprime_pairs}{checked},
\mword{Erdos249257/GcdMomentCalculus.lean}{349}{tsum_pos_coprime_inv_mersenne_eq_one}{checked},
\mword{Erdos249257/GcdMomentCalculus.lean}{474}{cylinderMass_split}{checked},
\mword{Erdos249257/TotientMahlerDefect.lean}{823}{paritySeparatedMatrix_det_ne_zero}{checked},
\mword{Erdos249257/TotientMahlerDefect.lean}{882}{exists_separatedMinorCertificate_totientAffineOddFamily}{checked},
\mword{Erdos249257/GenericTailOrbitRigidity.lean}{426}{binaryCoeffSeries_rational_iff_exists_temperedBinaryOrbit}{checked},
\mword{Erdos249257/GapFareyBound.lean}{51}{farey_gap}{checked},
\mword{Erdos249257/GapFareyBound.lean}{176}{gap_check_window_1_240_le_79639646646701375323355774875831053}{checked},
\mword{Erdos249257/CertificateKernel.lean}{15986}{tsum_totient_div_pow_two_ne_int_div_of_totient_gap_certificate}{checked},
\mword{Erdos249257/GapFareyBound.lean}{88}{gap_check_window_1_120_le_248672326362367909}{checked},
\mword{Erdos249257/CertificateKernel.lean}{18572}{tsum_visible_coprime_pairs_ne_int_div_of_den_le_39819823323350687661677887437915526}{checked},
\mword{Erdos249257/FreshPrimeDeficitDecomposition.lean}{171}{diagonalHeightIncrement_eq_old_add_endpointFreshDeficits}{diagonal split},
\mword{Erdos249257/TotientActualLcmShortKill.lean}{18}{lcmDiagonalArithmeticKill_two_pow_four}{the \(t=2^4\), length-\(23\) kill},
\mword{Erdos249257/TotientActualLcmShortKill.lean}{27}{lcmDiagonalArithmeticKill_two_pow_six}{the \(t=2^6\), length-\(93\) kill},
\mword{Erdos249257/TotientActualLcmShortKill.lean}{35}{actualLcmTailOrbit_four_and_six_notMem_int}{actual-orbit non-integrality at both heights},
\mword{Erdos249257/TotientActualLcmShortKill.lean}{54}{powerTwoActualLcmShortArithmeticKillSupply_through_six}{but explicitly stops at that finite prefix},
\mword{Erdos249257/TotientActualLcmOrbitSeparation.lean}{141}{abs_actualLcmTailOrbit_sub_rawApprox_lt}{recorded here},
\mword{Erdos249257/TotientActualLcmOrbitSeparation.lean}{179}{actualLcmRawApprox_eq_half_div_fourPow}{by the normalization identity},
\mword{Erdos249257/TotientActualLcmOrbitSeparation.lean}{208}{halfWordBandAt_of_rawApprox_integerSeparation}{through the exact threshold lemma},
\mword{Erdos249257/TotientActualLcmOrbitSeparation.lean}{254}{PowerTwoActualLcmOrbitSeparationSupply}{the cofinal separation hypothesis},
\mword{Erdos249257/TotientActualLcmOrbitSeparation.lean}{305}{irrational_totientSeries_of_actualLcmOrbitSeparationSupply}{the conditional endpoint theorem},
\mword{Erdos249257/TotientActualLcmTopEdgeStaircase.lean}{251}{not_actualLcmTerminalDyadicStaircase_of_room}{the staircase is impossible under the stated room hypothesis},
\mword{Erdos249257/TotientActualLcmTopEdgeStaircase.lean}{297}{windowDiscrepancy_emod_two_pow_eq_terminal}{terminal-remainder identity},
\mword{Erdos249257/TropicalCurvatureCarry.lean}{137}{fixedPrecisionTropicalNoGo}{checked},
\lword{Erdos249/TotientStrictPrimeEscape.lean}{157}{DTWNaturalPrimeTailOrbitStrictGap}{asks},
\lword{Erdos249/TotientStrictPrimeEscape.lean}{485}{naturalPivotPointEscape_of_naturalPrimeTailOrbitStrictGap}{a positive adaptive truncation budget},
\lword{Erdos249/TotientStrictPrimeEscape.lean}{524}{irrational_totient_series_of_naturalPrimeTailOrbitStrictGap}{$\Stot$ is irrational},
\href{https://github.com/wcook04/plectis-erdos/blob/99f4bf47422abbd8757cbb22b50ba079d764d3a7/ErdosProblems/Erdos249/TotientStrictPrimeEscape.lean#L35}{density hypothesis},
\href{https://github.com/wcook04/plectis-erdos/blob/99f4bf47422abbd8757cbb22b50ba079d764d3a7/ErdosProblems/Erdos249/TotientStrictPrimeEscape.lean#L25}{block-gap hypothesis},
\href{https://github.com/wcook04/plectis-erdos/blob/99f4bf47422abbd8757cbb22b50ba079d764d3a7/ErdosProblems/Erdos249/TotientStrictPrimeEscape.lean#L45}{density-to-gap implication},
\href{https://github.com/wcook04/plectis-erdos/blob/99f4bf47422abbd8757cbb22b50ba079d764d3a7/ErdosProblems/Erdos249/TotientStrictPrimeEscape.lean#L167}{exact squaring recurrence},
\href{https://github.com/wcook04/plectis-erdos/blob/99f4bf47422abbd8757cbb22b50ba079d764d3a7/ErdosProblems/Erdos249/TotientStrictPrimeEscape.lean#L191}{iterate formula},
\href{https://github.com/wcook04/plectis-erdos/blob/99f4bf47422abbd8757cbb22b50ba079d764d3a7/ErdosProblems/Erdos249/TotientStrictPrimeEscape.lean#L201}{initial-phase normal form},
\href{https://github.com/wcook04/plectis-erdos/blob/99f4bf47422abbd8757cbb22b50ba079d764d3a7/ErdosProblems/Erdos249/TotientStrictPrimeEscape.lean#L278}{initial-phase equivalence},
\href{https://github.com/wcook04/plectis-erdos/blob/99f4bf47422abbd8757cbb22b50ba079d764d3a7/ErdosProblems/Erdos249/TotientStrictPrimeEscape.lean#L322}{dyadic-root obstruction},
\href{https://github.com/wcook04/plectis-erdos/blob/99f4bf47422abbd8757cbb22b50ba079d764d3a7/ErdosProblems/Erdos249/TotientStrictPrimeEscape.lean#L465}{prime-alignment implication},
\lword{Erdos249/FiniteEulerSieve.lean}{28}{muSqEulerFactor_one}{$s=1$},
\lword{Erdos249/FiniteEulerSieve.lean}{36}{muSqEulerFactor_two}{$s=2$},
\lword{Erdos249/FiniteEulerSieve.lean}{44}{sigmaPrimePow_firstDiff}{first difference},
\lword{Erdos249/FiniteEulerSieve.lean}{51}{sigmaPrimePow_secondDiff}{second difference},
\lword{Erdos249/PrimeRayCyclotomicCurvature.lean}{22}{checkerboard}{mixed difference},
\lword{Erdos249/PrimeRayCyclotomicCurvature.lean}{26}{checkerboard_separable}{separable case},
\lword{Erdos249/PrimeRayCyclotomicCurvature.lean}{32}{checkerboard_unique}{fixed-stencil uniqueness},
\lword{Erdos249/PrimeRayCyclotomicCurvature.lean}{45}{fourPoint_layer_identity}{a four-point divisor-layer identity},
\lword{Erdos249/PrimeRayCyclotomicCurvature.lean}{62}{residueClass_step_has_centred_completion}{a centred representative for a residue class modulo an even modulus},
\lword{Erdos249/PrimeRayCyclotomicCurvature.lean}{95}{PrimeRayLayerSupply}{Layer},
\lword{Erdos249/PrimeRayCyclotomicCurvature.lean}{102}{BoundedDegreeOrderConsumer}{BoundedOrder},
\lword{Erdos249/PrimeRayCyclotomicCurvature.lean}{151}{FinitePrimeSupportEscape}{FinitePrimeEscape},
\href{https://github.com/wcook04/plectis-erdos/blob/99f4bf47422abbd8757cbb22b50ba079d764d3a7/Erdos249257/CarrySurvivorExtinction.lean#L515}{definition},
\mword{Erdos249257/TotientTailPeriodKiller.lean}{57}{totientTail}{definition},
\mword{Erdos249257/LcmConeFlatness.lean}{357}{not_irrational_of_tail_diff_mem_int}{checked},
\mword{Erdos249257/DiagonalPincerCertificatesT64.lean}{1933}{diagonalPincerCertificateScalesThroughT64}{list},
\mword{Erdos249257/DiagonalPincerCertificatesT64.lean}{1967}{certifiedKill_diagonal_all_imported_through_t64}{verification},
\lword{Skip/LadderT67.lean}{71264}{exists_diagonalKill_le_82}{checked},
\lword{Erdos249/CyclotomicAnchoredKill.lean}{3231}{CyclotomicAnchoredKillSupply}{CyclotomicAnchoredKillSupply},
\lword{Erdos249/CyclotomicAnchoredKill.lean}{3327}{binaryCyclotomicAnchoredKillSupply_iff_irrational}{the checked equivalence},
\lword{Erdos249/CyclotomicAnchoredKill.lean}{1706}{exists_unbounded_binaryCyclotomicSupport_with_periodLock_of_not_irrational}{checked conditional support/period no-go},
\lword{Erdos249/PrimeRayCyclotomicCurvature.lean}{158}{UnboundedPrimeDivisorSupply}{unbounded prime divisors},
\lword{Erdos249/CyclotomicAnchoredKill.lean}{1693}{binaryCyclotomicLayer_unboundedPrimeDivisorSupply}{the checked binary-layer theorem},
\lword{Erdos249/PeriodMultipleEscape.lean}{441}{ApFullDepthEscape}{ApFullDepthEscape},
\mword{Erdos249257/FirstHarmonicPivot.lean}{299}{mem_pivotFiber_iff_exists_supplierPrime}{membership},
\mword{Erdos249257/FirstHarmonicPivot.lean}{367}{pivotBaseOfPrime_injective_on_supplierPrimes}{injectivity},
\mword{Erdos249257/FirstHarmonicPivot.lean}{384}{image_pivotSupplierPrimes_eq_pivotFiber}{image equality},
\mword{Erdos249257/FirstHarmonicPivot.lean}{394}{supplierPrime_not_globally_isolated_counterexample}{checked counterexample},
\mword{Erdos249257/FirstHarmonicPivot.lean}{514}{windowFirstExp_sum_eq_pivot_decomposition}{checked},
\mword{Erdos249257/FirstHarmonicPivot.lean}{577}{irrational_totient_series_of_pivotResidualDecorrelation}{the conditional irrationality theorem},
\mword{Erdos249257/FirstHarmonicPivot.lean}{543}{PivotBudgetAt}{the budget condition},
\mword{Erdos249257/FirstHarmonicPivot.lean}{549}{first_harmonic_gap_of_pivotBudgetAt}{the first-harmonic gap},
\mword{Erdos249257/TotientKernelReduction.lean}{60}{totient_pow_mul_eq}{checked},
\mword{Erdos249257/TotientKernelReduction.lean}{117}{totient_pow_mul_affine_gcd_cross_eq}{checked}.

\begin{thebibliography}{9}\small
\bibitem{allouche-shallit} J.-P.~Allouche and J.~Shallit,
\emph{The ring of $k$-regular sequences}, Theoret.\ Comput.\ Sci.\
\textbf{98} (1992), 163--197. doi:10.1016/0304-3975(92)90001-V.
\bibitem{coons} M.~Coons, \emph{(Non)Automaticity of number theoretic
functions}, J.\ Th\'eor.\ Nombres Bordeaux \textbf{22} (2010), 339--352,
Theorem~3.2. doi:10.5802/jtnb.718.
\bibitem{martin-phi-inequalities} G.~Martin, \emph{Simultaneous inequalities
among values of the Euler phi-function}, 2006, arXiv:math/0603053,
Theorem~1 and Corollary~4.
\bibitem{wong2015} E.~Wong, answer to
\href{https://math.stackexchange.com/questions/1210886}{question 1210886},
Mathematics Stack Exchange, 29 March 2015.
\bibitem{erdosgraham1980} P.~Erd\H{o}s and R.~L.~Graham,
\emph{Old and New Problems and Results in Combinatorial Number Theory},
1980, p.~61.
\bibitem{fan2026totient} S.~Fan, comment on Erd\H{o}s Problem \#249,
16 May 2026, 19:01, \href{https://www.erdosproblems.com/forum/thread/249}
{Erd\H{o}s Problems discussion thread}.
\end{thebibliography}

\companionpapercontext

\end{document}
